\section{Canonical integration by parts with endpoint constants}
\label{poly:normalform}

Let $\mathcal A=\Q[u_0,u_1,u_2,\ldots]$ and
\[
 \cD=\sum_{r\ge0}u_{r+1}\frac{\partial}{\partial u_r},\qquad
 \wt(u_r)=r+1.
\]
Each polynomial uses finitely many variables. The weight-$w$ subspace
$\mathcal A_w$ has dimension $p(w)$, the partition number.
Let $\mathcal B_w$ consist of $u_0^w$ and all weight-$w$ monomials
whose highest positive-index variable occurs at least twice.
For $w=0$ use $\mathcal B_0=\{1\}$.

\begin{theorem}[A triangular normal form at every weight]
\label{poly:IBP}
For every $w\ge2$ there is a direct sum
\begin{equation}\label{poly:directsum}
 \mathcal A_w=\cD\mathcal A_{w-1}\oplus
                      \operatorname{span}_{\Q}\mathcal B_w.
\end{equation}
Every $P\in\mathcal A_w$ therefore has a unique decomposition
$P=\cD Q+N$ with $Q\in\mathcal A_{w-1}$ and $N$ in the displayed normal space. In particular
\begin{equation}\label{poly:dimension}
 |\mathcal B_w|=p(w)-p(w-1)\quad(w\ge2).
\end{equation}
The quotient has dimension one at weights zero and one, so its
Hilbert series is
\begin{equation}\label{poly:Hilbert}
 t+\prod_{j\ge2}(1-t^j)^{-1}.
\end{equation}
\end{theorem}

\begin{proof}
Suppose a monomial has highest variable $u_r$, $r\ge1$, occurring
exactly once. Write it as $u_r u_{r-1}^a V$, where $V$ involves
only $u_0,\ldots,u_{r-2}$. Then
\begin{equation}\label{poly:IBPstep}
 u_r u_{r-1}^aV=
 \cD\left(\frac{u_{r-1}^{a+1}V}{a+1}\right)
 -\frac{u_{r-1}^{a+1}\cD V}{a+1}.
\end{equation}
Every monomial in the last term has strictly smaller highest index.
Iterating therefore terminates and proves that the two spaces span.

For directness, let $Q$ be a nonconstant polynomial with highest
variable $u_m$. The highest-index terms of $\cD Q$ are
$u_{m+1}\partial Q/\partial u_m$, all with their highest variable
occurring once. They cannot occur in a polynomial in the normal space.
Since the characteristic is zero, $\partial Q/\partial u_m\ne0$.
This also proves that $\cD$ is injective on every positive-weight
subspace. Taking dimensions proves \eqref{poly:dimension}.
The weight-zero kernel consists of constants, which explains the extra
$t$ in \eqref{poly:Hilbert}.
\end{proof}

\begin{center}
\small
\begin{tabular}{ccl}
\toprule
Weight & Dimension & A canonical basis, with $u_r=\psi^{(r)}$\\
\midrule
2&1&$u_0^2$\\
3&1&$u_0^3$\\
4&2&$u_0^4,\ u_1^2$\\
5&2&$u_0^5,\ u_0u_1^2$\\
6&4&$u_0^6,\ u_0^2u_1^2,\ u_1^3,\ u_2^2$\\
7&4&$u_0^7,\ u_0^3u_1^2,\ u_0u_1^3,\ u_0u_2^2$\\
8&7&$u_0^8,\ u_0^4u_1^2,\ u_0^2u_1^3,\ u_1^4,
 u_0^2u_2^2,\ u_1u_2^2,\ u_3^2$\\
\bottomrule
\end{tabular}
\end{center}

\subsection{The endpoint map is part of the theorem}

For a polynomial $Q$, define
\begin{equation}\label{poly:endpoint}
 E(Q)=f_Q(1)-[x^0]f_Q(x).
\end{equation}
Both terms are finite polynomial calculations. The values at one are
$\psi(1)=-\gamma$ and
$\psi^{(r)}(1)=(-1)^{r+1}r!\zeta(r+1)$ for $r\ge1$.
The constant at zero is obtained from \eqref{poly:localpsi} with $n=0$;
it is not obtained merely by multiplying the individual constant terms.

\begin{corollary}[Exact finite-part integration by parts]
\label{poly:endpointcor}
For the normal decomposition $P=\cD Q+N$,
\begin{equation}\label{poly:IBPintegral}
 \boxed{\cI(P)=E(Q)+\cI(N).}
\end{equation}
For a common shift $a$ with $\re a>0$, the corresponding ordinary
identity is
\begin{equation}\label{poly:IBPshift}
 \int_0^1f_P(a+x)\dd x
 =f_Q(a+1)-f_Q(a)+\int_0^1f_N(a+x)\dd x.
\end{equation}
\end{corollary}
\begin{proof}
On $[\varepsilon,1]$ the integral of $f_{\cD Q}$ is
$f_Q(1)-f_Q(\varepsilon)$. Extracting its original-coordinate
finite part retains precisely $f_Q(1)-[x^0]f_Q(x)$.
The shifted identity is the ordinary fundamental theorem of calculus.
\end{proof}

For example, $\cD(u_0u_1)=u_1^2+u_0u_2$, but
\begin{equation}\label{poly:boundaryexample}
 E(u_0u_1)=-\zeta(3),\qquad
 \cI(u_1^2)+\cI(u_0u_2)=-\zeta(3).
\end{equation}
Likewise $E(u_1u_2)=4\zeta(5)$. These agree with all five bilinear
specializations above. Dropping total derivatives would lose these
constants and give a false relation.

The dimension theorem answers the stated formal classification problem
before imposing arithmetic identities among zeta coordinates.
Its numerical image may have smaller dimension. For instance the
$u_2^2$ coordinate at weight six already has the elementary zeta-derivative
evaluation \eqref{poly:square}. Neither the formal normal form nor a
failure to find further numerical relations proves arithmetic independence.
